% \subsection{Recent Methods to Achieve $w$-event LDP}
\subsection{\texorpdfstring{Recent Methods to Achieve $w$-event LDP}{Recent Methods to Achieve w-event LDP}}
\label{ss:competitors}

State-of-the-art $w$-event LDP methods~\cite{Ren2022} focus on releasing
statistical histograms at each timestamp.
These methods can be categorized into the following four types based on their
privacy budget allocation strategies.
\begin{itemize}
\item {LDP Budget Uniform (LBU)} applies a fixed $\epsilon/w$ budget to each
timestamp in the window.
However, as window size $w$ increases, the allocated budget per timestamp
becomes vanishingly small, introducing excessive noise and severely impairing
utility.

\item{LDP Sampling (LSP)} allocates the entire privacy budget $\epsilon$ to a
single timestamp within the sliding window, providing high accuracy at that
point while reusing its result to approximate the data for all other timestamps.
This method performs well when the stream is stable but fails to adapt when the
stream fluctuates.
The approximation errors can become very large if subsequent stream data differ
significantly from previous releases.

\item{LDP Budget Distribution (LBD)} and \textbf{LDP Budget Absorption
  (LBA)} adopt dynamic budget allocation techniques, drawing inspiration from
$w$-event CDP~\cite{Kellaris2014}.
Both methods consist of two sub-mechanisms: private dissimilarity estimation and
private strategy determination.
The private dissimilarity estimation sub-mechanism computes the private
dissimilarity between the current true statistics and the previous release based
on perturbed user data collected via an FO with a fixed budget
$\epsilon_{t,1}=\epsilon/(2w)$.
The private strategy determination sub-mechanism decides whether to publish new
statistics or reuse previous values by comparing the estimated dissimilarity
with a potential publication error determined by the publication budget
$\epsilon_{t,2}$.
The allocation of $\epsilon_{t,2}$ differs between LBD and LBA.
In LBD, the budget is distributed across timestamps requiring data publication
in an exponentially decaying manner, and budget spent at timestamps outside the
current window is reclaimed for reuse.
In contrast, LBA first allocates the budget uniformly and then absorbs it at the
timestamps that use approximation.

% However, adaptive mechanisms such as LBD and LBA, although intended
% to flexibly
% allocate privacy budgets according to data variation, have been reported to
% exhibit even lower efficiency in practice~\cite{Schaler2024},
% failing to fully
% exploit the benefits of dynamic budget allocation.
% In particular, the estimated data produced by sub-mechanism $M_{i,1}$, which
% consumes half of the privacy budget, is discarded immediately after computing
% the approximate error.
% This significantly reduces the effective budget available for data
% publication.
% Moreover, these strategies rely solely on current differences and
% the remaining
% budget, without considering correlations across the sliding window, which are
% essential under the $w$-event LDP model.

\item{Population Division Extensions.}
To improve the utility of LDP mechanisms, population division has emerged as a
promising strategy~\cite{Kulkarni2019,Wang2017l,Ren2022}.
Instead of splitting the privacy budget across timestamps, this paradigm
partitions the user population, assigning each user to report at a specific
timestamp using the entire privacy budget.
This reduces estimation variance via the amplification-by-subsampling effect and
enables more accurate data release within each reporting round.
Based on this idea, LBU, LBD, and LBA can be extended to LPU (Population-based
Uniform), LPD (Population-based Distribution), and LPA (Population-based
Absorption), respectively.
Notably, LSP intrinsically follows this paradigm.
Population division is a general augmentation strategy that can be applied to
any FO-based mechanism, including our proposed MTSP-LDP.
While population division improves utility within each reporting round, it does
not capture temporal correlations across timestamps, which are critical in
streaming scenarios.
\end{itemize}
% \bullethdr{Limitations of adaptive methods.}
\paragraph{Remarks}
Adaptive methods such as LBD and LBA, while designed to flexibly allocate
privacy budgets based on data variations, have been shown to exhibit even lower
utility in practice\cite{Schaler2024}, failing to leverage the potential
advantages of dynamic privacy budget allocation.
Specifically, the private dissimilarity estimation sub-mechanism consumes half
of the total privacy budget to collect perturbed data.
This data is then discarded after estimating dissimilarity, halving the budget
available for data publication.
Additionally, these methods rely solely on single-timestamp information,
neglecting the rich temporal correlations within the sliding window.
Consequently, this strategy is insufficient, as the $w$-event LDP constraint
requires the total budget $\epsilon$ to be allocated reasonably across all $w$
timestamps, not just based on single-timestamp changes and the instantaneous
budget cap.
